\section{Selection vs combination --- why online aggregation can pick, not blend}

\subsection{The idea}

We combine several forecasts into one. Two ways exist.
Combination: average the members with fixed or slow weights.
Selection: pick the single best member and follow it.
They sound similar. They are not.
Combination reduces variance. Selection does not.
This note explains when an ``adaptive weights'' method quietly
turns into selection --- and why that cost us MAE.

\subsection{Online aggregation and exponential weights}

Online aggregation updates the member weights every day.
It learns from how each member scored so far.
The classic recipe is \emph{exponential weights}.
Weight of member $m$ after day $D$:
\[
w_m \propto \exp\!\big(-\eta \, L_m\big),
\]
where $L_m$ is member $m$'s cumulative loss.
$\eta$ is the learning rate. A worse member gets a smaller weight.
This is the ``Hedge'' or exponentially weighted average forecaster.

\emph{Bernstein Online Aggregation} (BOA) is a modern refinement.
It adds a second-order term (the variance of the losses).
That lets it use a bigger learning rate safely.
It reacts faster and has a tighter theory guarantee.

\subsection{The objective is regret vs the best expert}

Here is the catch. What does BOA promise?
It bounds \emph{regret}: your total loss minus the best single
member's total loss. ``Expert'' is the theory word for member.
So the target is: do almost as well as the best single expert.

Read that again. The benchmark is the best \emph{one} member.
So the optimum BOA chases is: find that member, put all weight there.
When one member is clearly best, BOA's weights collapse onto it.
That is selection. The averaging is gone.

\subsection{Why selection forfeits variance reduction}

Combination wins because members make different mistakes.
Member A is $+2$ too high, member B is $-2$ too low.
Their average is perfect. The errors cancel.
This is variance reduction. It needs at least two members with weight.

Selection keeps only one member. Its errors have nobody to cancel.
So selection throws away the exact thing that makes a blend good.
A regret-vs-best-expert objective is built to select.
That is the wrong objective when diversity is your edge.

\subsection{Worked example: our numbers}

We tested BOA against our inverse-CRPS blend on the price product.
Same four members: LGBM, LEAR, Chronos, TFT.

\begin{tabular}{lcc}
\hline
Weighting scheme & MAE & Weight on LGBM \\
\hline
Inverse-CRPS (shipped) & 16.89 & shared across 4 \\
BOA online weights & 18.41 & 99.6\% \\
\hline
\end{tabular}

\vspace{0.5em}
BOA put 99.6\% of the weight on LGBM.
It selected the single best member and dropped the rest.
Result: 18.41 MAE --- 1.52 worse than the blend.
It even lost to LGBM alone, because the last 0.4\% weight
on weaker members added noise without enough averaging to pay for it.
The inverse-CRPS blend keeps real weight on all four.
The diversity does the work.

\subsection{When BOA IS the right tool}

BOA is not bad. It is a mismatch here. It shines when:
\begin{itemize}
\item You have many experts, and they are highly correlated.
\item One expert is clearly best, but you do not know which in advance.
\item The best expert may change over time (regime shifts).
\end{itemize}
Then selection IS near-optimal, and BOA's regret bound is real value.
It tracks the winner without you hand-picking it.
Our case is the opposite: few members, deliberately diverse,
each contributing. So a fixed convex average beats adaptive selection.

\subsection{Interview line}

``I tested Bernstein Online Aggregation against my inverse-CRPS
blend. BOA piled 99.6\% of the weight onto the best single model
and scored 18.41 MAE versus the blend's 16.89. The lesson: BOA
minimizes regret against the best expert, so it converges to
selection. When your edge is error diversity, selection forfeits
the variance reduction. Adaptive weights are the right tool only
when many correlated experts hide one clear winner.''
